%----------------------------------------------------------------------------
\chapter{The TDK lifecycle}\label{app:lifecycle}
%----------------------------------------------------------------------------
Plan the phases in \cref{tab:lifecycle} backwards from the official dates in \cref{app:tdk}: the upload date is fixed, so every step before it is that date minus the time the rest needs. Each step has an owner, a condition that says when it is done, and a link to the guidance for it.
We track every student's progress in the same steps, so at the weekly consultation (\cref{app:consult}) we look at the row you are in and the row you are heading for. Some steps are marked \emph{a hill climb}: you repeat them, scoring the work, fixing its weakest point and scoring again, until the row's condition says stop.
Set \texttt{\textbackslash lintphase} in \texttt{thesis.tex} to your phase's name in lower case, from \texttt{topic} to \texttt{submit}; the lint's checks come due by phase (\cref{app:code:paperci}). After the upload, leave it at \texttt{submit}: the paper is submitted and needs no further checks.
Inspect the reviewer PDF every week. It is the version reviewers will read.

Three steps are \anchor{}s: the topic manifest, the spine and the full draft. An anchor is an artifact we approve before the next phase, and everything after it is measured against it. The topic manifest and the spine are one page each; we approve the full draft at a revision recorded in the consultation summary. Each anchor narrows the topic: before the manifest it is a direction, after the spine it is a fixed set of claims with fixed names, and after the full draft it is a paper that only gets polished. Nothing stops at an anchor, but work that contradicts one is drift, and drift is the most expensive mistake in the process, because everything built on the old anchor has to be redone. If an anchor turns out to be wrong, say so at the next consultation and we change it together; a claim changed silently costs weeks when a reviewer finds the seam.

Red rows mark anchors. Yellow rows marked \official{} are the faculty's steps: their portal dates (\cref{app:tdk:dates}) are fixed, and missing one ends that semester's TDK. Green rows marked \department{} are the department's sessions, where every TDK student of the department comes and the last problems get fixed in company; we organise them.

\newcommand{\lcphase}[2]{\midrule\multicolumn{4}{l}{\textbf{#1}\quad\textit{#2}}\\\midrule}
% Support cells are links with a name, not references with a number: \lcref{label}{name}.
\newcommand{\lcref}[2]{\hyperref[#1]{#2}}
{\small
\begin{xltabular}{\linewidth}{>{\raggedright\arraybackslash}p{0.2\linewidth} >{\raggedright\arraybackslash}p{0.07\linewidth} >{\raggedright\arraybackslash}X >{\raggedright\arraybackslash}p{0.19\linewidth}}
  \caption{TDK steps, completion conditions and guidance. Owners are you, the advisors (we), or both.}\label{tab:lifecycle}\\
  \toprule
  \textbf{Step} & \textbf{Owner} & \textbf{Done when} & \textbf{Support} \\
  \endfirsthead
  \toprule
  \textbf{Step} & \textbf{Owner} & \textbf{Done when} & \textbf{Support} \\
  \endhead
  \bottomrule
  \endlastfoot
  \lcphase{0. Topic}{before the semester}
  \rowcolor{anchorrow}
  Topic manifest \anchor & both & One page, agreed with us: the problem, why it matters, the expected contribution & the TDK topic page \\
  \lcphase{1. Setup}{first week}
  Identifiers & you & We have your Neptun code, GitHub username and e-mail & \lcref{app:start:ids}{Share your identifiers} \\
  Overleaf and paper repository & we & You can edit in Overleaf and push to the synced GitHub repository & \lcref{app:start:access}{Accounts and access} \\
  Repositories & you & Notes repository with \texttt{paper/} and \texttt{code/} as submodules; both advisors added to all three & \lcref{app:start:access}{Accounts and access} \\
  Zotero & both & You are in the group, the connector is installed, \texttt{bib/zotero.bib} is linked to the group & \lcref{app:start:access}{Accounts and access} \\
  Calendar and weekly slot & you & Neptun timetable synced, calendars shared, a recurring Teams event & \lcref{app:start:calendar}{Calendar}, \lcref{app:consult}{Weekly consultations} \\
  Virtual machine & both & You asked for the VM and can reach it over SSH & \lcref{app:start:env}{Development environment} \\
  Background material & we & Our papers and notes on the topic are in your notes repository & --- \\
  \lcphase{2. Research}{first half of September}
  Problem & you & The problem is stated, with a running example & \lcref{ch:intro}{Introduction} \\
  Literature review and gap & you & The search is logged in one notes file, with its search strings, databases and dates; the sources are in Zotero; the positioning figure and the gap matrix are filled; the gap is one sentence & \lcref{fig:positioning}{positioning figure}, \lcref{tab:related}{gap matrix}, \lcref{app:writing:refs}{The reference list} \\
  Research questions and contributions & you & RQ1--RQ3 and C1--C3 are written as checkable claims & \lcref{ch:intro}{Introduction}, \lcref{ch:evaluation}{Evaluation} \\
  Threats and requirements & you & Threat model, assumptions and requirements are listed & \lcref{ch:approach}{Approach} \\
  Leading figure draft & you & The figure explains the work without the text & \lcref{app:writing:figure}{The leading figure} \\
  Proof of concept & you & The core idea runs end to end on the running example; shortcuts are fine & \lcref{app:code:github}{GitHub} \\
  \lcphase{3. Plan}{second half of September}
  Running example and validation plan & you & The example is named and appears in every chapter; how each claim will be checked is written down & \lcref{app:writing:spine}{Plan before you write} \\
  \rowcolor{anchorrow}
  Spine \anchor & both & One page of claims, terms and notation, agreed with us & \lcref{app:writing:spine}{Plan before you write} \\
  Skeleton mind map & you & Each chapter has paragraph topic sentences and evidence for each claim & \lcref{fig:mindmap}{the example skeleton} \\
  Premortem & both & The reviewer panel has read the spine; likely objections for every review-sheet criterion are answered in the plan, and your decisions are in the consultation summary & \lcref{app:lifecycle:reviewers}{Reviewer passes}, \lcref{app:tdk:sheet}{The review sheet} \\
  \rowcolor{official}
  Registration with the abstract \official & you & Registered on the portal with the Hungarian and English abstract, by 30 September & \lcref{app:tdk:dates}{Official dates}, \lcref{app:writing:abstract}{The abstract} \\
  Complete, testable implementation & you & By 30 September every component exists, the evaluation workloads run, and a test suite runs & \lcref{ch:implementation}{Implementation} \\
  \lcphase{4. Draft}{October}
  Outline check & we & The skeleton fits the template's chapters and the review sheet & the chapter todos \\
  Proofs (a hill climb) & you & Every claim has its proof, mechanised proof or measurement & \lcref{ch:approach}{Approach} \\
  Tested implementation (a hill climb) & you & All tests pass; \texttt{just reproduce} regenerates every result; the release is ready for Zenodo & \lcref{app:code:github}{GitHub} \\
  Scalability and robustness & you & Measurements beyond small examples, with setup and repetitions, and a robustness analysis & \lcref{fig:guide:scalability}{scalability figure}, \lcref{tab:guide:robustness}{robustness table} \\
  Panel pass & you & Every relevant persona has read the draft; blockers are fixed or acknowledged in the limitations before we receive it & \lcref{app:lifecycle:reviewers}{Reviewer passes} \\
  \rowcolor{anchorrow}
  Full draft \anchor & both & Every chapter is written and in our hands ten days before the deadline; we approve it at a revision recorded in the consultation summary & --- \\
  \lcphase{5. Polish}{22--28 October}
  Advisor review & both & Our comments are resolved and the todo list is empty & \lcref{app:latex:todo}{Todo notes} \\
  Review sheet hill climb & you & You have scored A--E, fixed the weakest criterion and repeated until no cheap fixes remain & \lcref{app:tdk:sheet}{The review sheet} \\
  Skim pass (a hill climb) & you & Three passes done, or fewer once we accept what remains; aim for the maximum on every criterion from the skim view alone & \lcref{app:lifecycle:reviewers}{Reviewer passes} \\
  Code archive & you & The code is on Zenodo, cited by its concept DOI, and the results tag is named in the paper & \lcref{app:zenodo}{Zenodo} \\
  \rowcolor{department}
  Writing session \department & both & You attended with pizza, closed the last todos and had an advisor check the reviewer PDF & \lcref{tab:polish}{Requirements} \\
  \lcphase{6. Submit}{29 October -- 1 November, or earlier once the paper is complete}
  Requirements & you & \texttt{\textbackslash lintphase} is \texttt{submit}, every row due by \texttt{submit} is met and \texttt{make check-submit} is clean & \lcref{tab:polish}{Requirements} \\
  \rowcolor{official}
  Paper upload \official & you & AI declaration filled in, main text 30--35 pages, the reviewer PDF uploaded under \emph{Dolgozataim} by 1 November & \lcref{app:tdk:form}{Form requirements} \\
  \lcphase{7. After the deadline}{2--16 November, while the reviewers read}
  Code polish & you & Within about two days of upload: README complete, Docker images build, results map filled, a new Zenodo version published; the results tag untouched. Later polish is allowed, but reviewers may already have looked & \lcref{app:zenodo:versions}{Release a new version} \\
  Demo & you & The demonstration is recorded by a script and published by CI on every push, with captions that carry the insight & \lcref{app:code:demo}{Demos} \\
  Presentation & you & Slides in the group's template, the leading figure second, one story in fifteen minutes, rehearsed about ten times & \lcref{app:lifecycle:presentation}{The presentation}, \href{https://github.com/ftsrg/cheat-sheets/wiki/TDK-Howto}{FTSRG TDK Howto} \\
  Reviews and questions & you & Every review question, doubt and remaining panel objection has an answer prepared, with backup slides where useful & \lcref{app:tdk:sheet}{The review sheet}, \lcref{app:lifecycle:reviewers}{Reviewer passes} \\
  \rowcolor{department}
  Presentation rehearsal \department & both & You rehearsed two days before the conference, practised answering the reviewers' questions and fixed the problems found & \lcref{app:lifecycle:presentation}{The presentation} \\
  \lcphase{8. Conference}{18 November}
  \rowcolor{official}
  Conference \official & you & You gave the presentation in your section & \lcref{app:lifecycle:presentation}{The presentation}, \lcref{app:tdk:dates}{Official dates} \\
  After the TDK & both & Decided together: the OTDK if nominated, a conference paper, or a thesis & \lcref{app:tdk:dates}{Placements} \\
\end{xltabular}
}

\section{Reviewer passes}\label{app:lifecycle:reviewers}
Two reviewers each spend an evening reading and scoring the paper against the review sheet (\cref{app:tdk:sheet}); the section jury, mostly from other fields, hears your presentation and reads little else. These are different readings, and a paper that satisfies you and us can still fail one of them. Prepare for each with a \emph{reviewer pass}: an agent in your paper repository reads the paper as a simulated reader with a narrow brief and reports what that reader would hold against it. A pass takes minutes, so run it early and repeat it. Personas can miss problems or invent them, so treat their output as objections to check, not as an assessment. \Cref{app:code:reviewers} says how to run them.

\paragraph{The skim pass.} The template aims to earn the maximum score from a reviewer who reads only the outline, figures, labelled items and coloured boxes, because a tired reviewer on the second evening reads exactly that. Test this with a \emph{skim view} made from the reviewer PDF, which paper CI builds on every push. It keeps the title, abstract, labelled items and boxes verbatim. Headings retain their page and length; figures retain their captions; tables retain their captions and a thumbnail. Paragraphs, equations, listings and plain lists become length markers, such as \enquote{paragraph, 7 lines}.

The skim reviewer reads only this view and the review sheet. It scores each criterion and, for any shortfall, names the item that should have made the point clear. For example, a guarantee may lack a row in the correctness table. During polish, repeat the pass and fix what it names. Aim for the maximum on every criterion, but stop after three passes, or earlier once we accept the remaining objections: past that point the fixes cost more than the points they recover. This is the review sheet hill climb for a reader who skims.

\paragraph{The panel.} \Cref{tab:reviewers} lists the expert personas. Their responsibilities are \emph{disjoint}: each concern has one owner and one place for an answer, so two personas never argue the same point and you never get the same objection twice in different words. A persona that notices a concern outside its remit hands it to the owner in one line, without arguing it.

Five rules define these responsibilities. The cryptographer and formal methods expert handle claims inside the formal model, and the formal methods expert owns the choice of abstraction: what the model leaves out and whether leaving it out is safe. The networking and cybersecurity expert owns the attacks the paper excludes. The mean mathematician checks pen-and-paper proofs; the formal methods expert checks tool-assisted proofs. The systems modelling expert checks system structure; the process modelling expert checks behaviour. The distributed systems expert checks which faults and scales are studied; the empiricist checks how they are measured. Every persona checks overclaiming in its own area. Only the champion looks for results stated more weakly than the evidence supports.

The first four personas read every paper. The others read papers that touch their field.
{\small
\begin{xltabular}{\linewidth}{>{\raggedright\arraybackslash}p{0.2\linewidth} X >{\raggedright\arraybackslash}p{0.22\linewidth}}
  \caption{Reviewer personas, their responsibilities and the papers they read. Each concern has one owner.}\label{tab:reviewers}\\
  \toprule \textbf{Persona} & \textbf{Owns} & \textbf{Reads} \\ \midrule \endfirsthead
  \toprule \textbf{Persona} & \textbf{Owns} & \textbf{Reads} \\ \midrule \endhead
  \bottomrule \endlastfoot
  Skim reviewer & Whether the skim view alone earns the maximum on every criterion & every paper \\
  Out-of-field reviewer & The problem and why it matters on the first page; every term explained before it is used; the leading figure without the text; the conclusion in plain words & every paper \\
  Reviewer two & Novelty against the closest neighbour; the gap matrix and the positioning figure; recent research in the literature; whether the topic is still open & every paper \\
  Champion & Results reported weaker than the evidence allows; contributions missing from the abstract, the leading figure or the conclusion & every paper \\
  Mean mathematician & Definitions and notation; theorem statements with every assumption named; each step of the pen-and-paper proofs & definitions, theorems, proofs \\
  Formal methods expert & The choice of abstraction: what the formal model leaves out and whether leaving it out is safe; whether each property says what the prose claims; vacuity, bounds and what the tool trusts; translations into the tool's input & model checkers, provers, formal semantics \\
  Cryptographer & The adversary in the model; the primitives and their assumptions; security properties as definitions, and the argument for them; composition & cryptography, credentials, privacy \\
  Systems modelling expert & The metamodel or system model and its constraints; a modelling language the paper defines; requirements traced to model elements; one name per element across figures and text & metamodels, DSLs, system and architecture models \\
  Distributed systems expert & Timing and fault models; safety and liveness; consistency and finality; whether the scalability and robustness scenarios exercise the fault model & ledgers, consensus, replication, clocks \\
  Networking and cybersecurity expert & Attacks the paper excludes: the network, keys and identities in practice, implementation and deployment, metadata; a reason for every exclusion & networks, smart contracts, keys \\
  Process modelling expert & Conformance to the notation; deadlock and proper completion; exception and timeout paths; the execution semantics assumed & BPMN, workflows, choreographies, sequence and state diagrams \\
  Empiricist & Setup, workloads, baselines, repetitions and statistics; precision and recall of a tool; threats to validity; whether \texttt{just reproduce} regenerates the figures & measurements, tool evaluations \\
\end{xltabular}
}

\paragraph{When.}\label{app:lifecycle:when-passes} Run the passes at four steps of \cref{tab:lifecycle}. First, use the panel for the spine's premortem: predict the objections to the claims and answer them in the plan, while changing a claim still costs a sentence rather than a chapter. Second, have every relevant persona read the draft, and fix each blocker or acknowledge it in the limitations before the full draft reaches us, so that our reading goes to what only we can see. Third, repeat the skim pass during polish. Fourth, after the deadline, turn the remaining panel objections into the reviewer questions you rehearse answers to.

\section{Requirements}\label{app:lifecycle:polish}
The words \emph{must}, \emph{we recommend} and \emph{you can} mean what \cref{app:here:words} says.

\Cref{tab:polish} lists the requirements, in the order they come due. A rule name marks a row the lint checks; \emph{reading} marks a row you check yourself. The lint's rows are checked by \texttt{make check-reviewer}, which runs every rule against the reviewer PDF (\cref{app:code:paperci}). The paper CI runs the same rules on every push, so by the last week they should have been green for a while: a red row in the last week is a surprise, and the last week has no room for surprises. The polish phase is for checking, not writing: if a row needs more than an hour, it belonged to the draft phase, and the full draft was not full. A row that restates a step of \cref{tab:lifecycle} only names the step and points there. Before upload, set \texttt{\textbackslash lintphase} to \texttt{submit}, run \texttt{make check-submit}, which runs the lint at that phase whatever \texttt{\textbackslash lintphase} says, and go through every row due by \texttt{submit}.
{\small
\begin{xltabular}{\linewidth}{X >{\raggedright\arraybackslash}p{0.12\linewidth} >{\raggedright\arraybackslash}p{0.15\linewidth} >{\raggedright\arraybackslash}p{0.1\linewidth}}
  \caption{Requirements, their source, what checks them and the phase they are due by. Rule names identify lint rules (\cref{app:code:paperci}); \emph{reading} means you check the row yourself.}\label{tab:polish}\\
  \toprule \textbf{Requirement} & \textbf{Source} & \textbf{Checked by} & \textbf{Due} \\ \midrule \endfirsthead
  \toprule \textbf{Requirement} & \textbf{Source} & \textbf{Checked by} & \textbf{Due} \\ \midrule \endhead
  \bottomrule \endlastfoot
  \textbf{Weekly consultation.} You come every week, even if little has changed (\cref{app:consult}). & Advisors & reading & every week \\
  \textbf{Consultation summary.} A short summary of what we agreed is saved as \texttt{notes/consultations/YYYY-MM-DD*.md}. & Advisors & reading & every week \\
  \textbf{Recorded call.} The consultation's Teams call is recorded and transcribed. & Advisors & reading & every week \\
  \textbf{Phase.} \texttt{\textbackslash lintphase} in \texttt{thesis.tex} names your phase in \cref{tab:lifecycle}. & Advisors & \texttt{frame} & every week \\
  \textbf{Reviewer PDF inspected.} You read this week's reviewer PDF; it is the version reviewers will read. & Advisors & reading & every week \\
  \textbf{Build log.} No undefined or multiply defined labels, no undefined citations, no overfull boxes, no missing fonts or files. & Advisors & \texttt{build}, \texttt{overfull} & every week \\
  \textbf{Cross-references.} Every reference uses \texttt{\textbackslash cref}; every figure, table, listing and equation is referenced from the text (\cref{app:latex:cref}). & Advisors & \texttt{figures}, \texttt{crefs} & every week \\
  \textbf{Own preamble.} \texttt{include/macros.tex} follows \cref{app:latex}. & Advisors & \texttt{preamble} & every week \\
  \textbf{Quotations.} Every verbatim quotation is in \texttt{\textbackslash enquote} and followed by a citation; paraphrases are cited too (\cref{app:writing:refs}). & Advisors & reading & every week \\
  \textbf{AI conversations kept.} Every AI conversation about the work is saved, none in a temporary or incognito chat (\cref{app:writing:ai}). & Advisors & reading & every week \\
  \textbf{AI answers verified.} You verified every AI answer you used; you are responsible for all AI-generated content. & Faculty: AI guideline & reading & every week \\
  \textbf{Essential work your own.} AI use is not disproportionately large and does not touch the essential parts of the work (\cref{app:writing:ai}). & Faculty: AI guideline & reading & every week \\
  \textbf{Topic manifest.} The anchor is approved (\cref{tab:lifecycle}, Topic manifest). & Advisors & reading & \texttt{topic} \\
  \textbf{Repositories.} The notes repository has \texttt{paper/} and \texttt{code/} as submodules; both advisors are added to all three (\cref{app:start:access}). & Advisors & reading & \texttt{setup} \\
  \textbf{Literature search log.} One notes file logs the search: its search strings, databases and dates. & Advisors & reading & \texttt{research} \\
  \textbf{Registration.} You registered on the portal with the Hungarian and English abstract by 30 September; missing it ends this semester's TDK (\cref{app:tdk:dates}). & Faculty: call & reading & \texttt{plan} \\
  \textbf{Spine.} One page of claims, terms and notation, approved by us (\cref{app:writing:spine}). & Advisors & reading & \texttt{plan} \\
  \textbf{Running example.} It is named with \texttt{\textbackslash exname}, introduced in \tplref{ch:intro} and shown in a box in every chapter (\cref{app:writing:spine}). & Advisors & \texttt{frame}, \texttt{essentials} & \texttt{plan} \\
  \textbf{Complete, testable implementation.} The step is done by 30 September (\cref{tab:lifecycle}, Complete, testable implementation). & Advisors & reading & \texttt{plan} \\
  \textbf{Outline.} Chapters are renamed where the template says so, no heading is a TODO; each chapter is one file in reading order, the introduction first and nothing after the conclusion. & Advisors & \texttt{outline}, \texttt{chapters}, \texttt{spine} (chapter order) & \texttt{draft} \\
  \textbf{Full draft.} The anchor is in our hands ten days before the deadline and approved (\cref{tab:lifecycle}, Full draft). & Advisors & reading & \texttt{draft} \\
  \textbf{Essentials.} Every essential has reached its highest level, written and used (\cref{app:writing:essentials}). & Advisors & \texttt{essentials} & \texttt{polish} \\
  \textbf{Every chapter written.} Chapters outside the essentials, such as background or implementation, have at least 150 words of their own, beyond the template skeleton or an outline. Every file containing text is included in the paper. & Advisors & \texttt{stubs}, \texttt{orphans} & \texttt{polish} \\
  \textbf{Figure sources.} Every figure is TikZ in the paper repository or an SVG included with \texttt{\textbackslash includesvg}; a raster image or screenshot has a written reason and is not called vector (\cref{app:writing:figures}). & Advisors & \texttt{figures}, reading & \texttt{polish} \\
  \textbf{Yours is visible.} The leading figure and every architecture-like figure colour your components as \cref{app:writing:figure} says. & Advisors & reading & \texttt{polish} \\
  \textbf{Axes and legends.} Every axis has a label with a unit; every legend entry names a thing the text also names. & Advisors & \texttt{axes}, reading & \texttt{polish} \\
  \textbf{Captions.} Every caption follows \cref{app:writing:captions}. & Advisors & \texttt{captions}, reading & \texttt{polish} \\
  \textbf{Figure text.} 10\,pt in print is the target, as the call says; 8\,pt is the floor we accept where 10\,pt does not fit. The lint flags only text below 8\,pt, so check text between 8 and 10\,pt yourself (\cref{app:writing:figures}). Colour pairs are distinguishable in greyscale. & Faculty: call & \texttt{pdf-smalltext}, reading & \texttt{polish} \\
  \textbf{Tables.} Tables fit the text width, numbers are right-aligned with the same precision in a column, units are in the header. & Advisors & \texttt{tables}, reading & \texttt{polish} \\
  \textbf{Spine drift.} Every term, symbol and claim matches the spine; one name per concept across chapters. & Advisors & reading & \texttt{polish} \\
  \textbf{English throughout.} No Hungarian outline sentence is left in a chapter; the whole paper, title and abstract included, passes a spell check; British or American spelling is consistent. & Advisors & \texttt{language}, \texttt{spelling} & \texttt{polish} \\
  \textbf{Citations.} Every claim from the literature is cited, the reference list had its repair pass (\cref{app:writing:refs}), and capitalisation is protected with braces (\cref{app:latex:cite}). & Advisors & \texttt{bib}, reading & \texttt{polish} \\
  \textbf{Numbered references.} The bibliography is numbered through \texttt{tdkieee.bst} (\cref{app:latex:cite}). & Advisors & reading & \texttt{polish} \\
  \textbf{Long reference list.} Above 50 entries, each one supports a claim or a fair comparison; there is no target to cut to (\cref{app:writing:refs}). & Advisors & reading & \texttt{polish} \\
  \textbf{No weak sources.} No encyclopedia, tutorial, vendor marketing or unreviewed post is cited as evidence (\cref{app:writing:refs}). & Advisors & reading & \texttt{polish} \\
  \textbf{Web sources.} A web source in the bibliography names its owner, title, URL and access date (\cref{app:latex:cite}). & Advisors & \texttt{bib} & \texttt{polish} \\
  \textbf{Abbreviations, notation, terms.} Every abbreviation, symbol and term is declared and used through \texttt{\textbackslash gls} (\cref{app:latex:acro}). & Advisors & \texttt{acronyms} & \texttt{polish} \\
  \textbf{Paper and layout.} A4 paper, single column (\cref{app:tdk:form}). & Faculty: call & \texttt{pdf-file}, reading & \texttt{polish} \\
  \textbf{Body font.} Times New Roman or another readable serif, at 12\,pt. & Faculty: call & \texttt{pdf-font} & \texttt{polish} \\
  \textbf{Line spacing.} Single. & Faculty: call & \texttt{pdf-spacing} & \texttt{polish} \\
  \textbf{Margins.} Between 1.5 and 3.5\,cm on every side. & Faculty: call & \texttt{pdf-margins} & \texttt{polish} \\
  \textbf{PDF size.} At most 15\,MB, our safeguard taken from the faculty's thesis guideline (\cref{app:code:paperci}). & Advisors & \texttt{pdf-file} & \texttt{polish} \\
  \textbf{PDF.} Bookmarks and metadata are correct, every link resolves, and the Zenodo DOI opens the record. & Advisors & \texttt{pdf-file}, \texttt{links}, reading & \texttt{polish} \\
  \textbf{Read-only passes.} Three readers can each follow one layer alone: figures and captions; running-example boxes; or the abstract and each chapter's first paragraph. Fill in any step missing from a layer. & Advisors & reading & \texttt{polish} \\
  \textbf{Reproducibility.} \texttt{docker run \dots{} reproduce} regenerates every figure and table from the results tag named in the paper (\cref{app:code:readme}). & Advisors & reading & \texttt{polish} \\
  \textbf{Code archive.} The code is on Zenodo, cited by its concept DOI, and the results tag is named in the paper (\cref{app:zenodo:cite}). & Advisors & \texttt{essentials}, \texttt{placeholders} & \texttt{polish} \\
  \textbf{Review sheet hill climb.} Finished: no cheap fix is left on any criterion (\cref{tab:lifecycle}, Review sheet hill climb; \cref{app:tdk:sheet}). & Advisors & reading & \texttt{polish} \\
  \textbf{Reviewer passes.} The skim pass is finished and a run of the panel raises no blocker (\cref{tab:lifecycle}, Skim pass; \cref{app:lifecycle:reviewers}). & Advisors & reading & \texttt{polish} \\
  \textbf{Page count.} The main text, from the first page of \tplref{ch:intro} to the last page of the conclusion (\path{content/07-conclusion.tex}), is 30--35 pages; the title page, the abstracts, the acknowledgements, the glossary lists, the bibliography, the AI declaration and the appendices do not count (\cref{app:tdk:form}). & Faculty: call; the counting rule: us & \texttt{pages}, \texttt{pdf-pages} & \texttt{submit} \\
  \textbf{Both abstracts.} The Kivonat and the Abstract are in the PDF (\cref{app:writing:abstract}). & Advisors & \texttt{essentials}, \texttt{pdf-parts} & \texttt{submit} \\
  \textbf{AI declaration.} It is filled in: one box marked and, if you used AI, each use with its tool, parts, share and prompt, and the overall share justified (\cref{app:writing:ai}). & Advisors & \texttt{aidecl}, \texttt{pdf-parts} & \texttt{submit} \\
  \textbf{AI citations.} Each AI answer is cited in the declaration table in the faculty's MLA form, never in the bibliography; a fact taken from it is cited to the primary source that verifies it (\cref{app:writing:ai}). & Advisors & reading & \texttt{submit} \\
  \textbf{Front and back matter.} Title page variables, both abstracts, the AI declaration, acknowledgements and appendices are complete. No faculty-template placeholder text, \enquote{TODO} or placeholder figure remains. & Advisors & \texttt{placeholders}, \texttt{stock}, \texttt{titlepage} & \texttt{submit} \\
  \textbf{Notes gone.} Every beat has become prose and every todo and advisor note is resolved; the burn-down counts are zero. & Advisors & \texttt{beats}, \texttt{markers} & \texttt{submit} \\
  \textbf{Upload.} The reviewer PDF is uploaded under \emph{Dolgozataim} by 1 November; missing it ends this semester's TDK (\cref{app:tdk:dates}). & Faculty: call & reading & \texttt{submit} \\
  \textbf{Code polish.} Released as a new Zenodo version within about two days of upload, the results tag untouched (\cref{tab:lifecycle}, Code polish; \cref{app:zenodo:versions}). & Advisors & reading & after the deadline \\
  \textbf{Demo.} The demonstration is recorded by a script and published by CI on every push (\cref{app:code:demo}). & Advisors & reading & after the deadline \\
  \textbf{Presentation.} You give it in your section on 18 November, in fifteen minutes plus five for questions; missing it ends this semester's TDK (\cref{app:tdk:dates}). & Faculty: call & reading & conference \\
\end{xltabular}
}

\section{The presentation}\label{app:lifecycle:presentation}
You give the presentation on conference day, in your section; \cref{app:tdk:dates} gives the date and the format. The Presentation, Reviews and questions, and Presentation rehearsal rows of \cref{tab:lifecycle} say how to prepare it.
\advisortodo{Link the group's slide template (Beamer and PowerPoint) here and in the Presentation row of \cref{tab:lifecycle}.}
\begin{itemize}
  \item \href{https://www.youtube.com/watch?v=Unzc731iCUY}{MIT: How to Speak}, for the presentation.
\end{itemize}
